\chapter*{}
\noindent{\textbf{RESUMO}}

\noindent{A disponibilidade de genomas de espécies e variedades de \textit{Coffea} permite comparar suas proteínas, o que exige localizar o gene correspondente em cada proteoma e adotar uma medida de similaridade explícita. Este trabalho apresenta um protótipo de ferramenta para a análise comparativa de proteínas de vias associadas ao sabor e ao aroma do café. A partir de uma sequência-âncora, o gene é localizado por BLASTP, com confirmação por melhores correspondências recíprocas, e as proteínas são alinhadas por uma implementação do algoritmo de Gotoh. Foram comparados os genes \textit{DXMT}, \textit{HCT} e \textit{SPS} em quatro proteomas. HCT e SPS foram localizados em todos, com identidade entre 92,56\% e 100\%; o DXMT, em dois, com 98,96\%. A confirmação recíproca evitou a escolha de um parálogo. A API e a interface web serão desenvolvidas em seguida.}

\vspace{\onelineskip}

% todas em letras minúsculas, separadas por ponto e vírgula (;)
\noindent{\textbf{Palavras-chave}: bioinformática; alinhamento de sequências; \textit{Coffea}; ortologia; similaridade de proteínas.}
